\chapter{Theory and Fundamentals of NeuroLearn}

This chapter presents the theoretical foundations and core principles behind the proposed AI-driven adaptive learning platform. The system combines cognitive learning theories, behavioral analytics, retrieval-augmented generation, accessibility-first design, and multi-modal content delivery to create a personalized educational experience. The chapter also discusses the hybrid learner profiling mechanism, dynamic modality adaptation, accessibility integration, and the gesture-based hardware module developed for inclusive interaction.

\section{Dynamic Learning Style Modeling}

The platform is built upon the VARK framework, which classifies learners into Visual, Auditory, Reading/Writing, and Kinesthetic modalities. Unlike traditional systems that treat learning styles as static traits, the proposed system models learner preference as a dynamic behavioral signal. The learner initially completes a VARK questionnaire, generating a baseline modality vector. Over time, real-time interaction behavior such as content replay frequency, time spent on simulations, use of audio explanations, and navigation patterns are continuously monitored. The final learner profile is computed using exponential decay fusion between the questionnaire profile and behavioral evidence, enabling the system to gradually shift from self-reported preference to observed interaction behavior.

\section{Multi-Modal Content Generation}
The adaptive content engine generates learning material according to the learner’s dominant modality. Visual learners receive animated diagrams, concept maps, and graphical explanations generated using tools such as Matplotlib, Graphviz, and Manim. Auditory learners receive AI-generated speech explanations and voice-guided walkthroughs using neural text-to-speech systems. Reading/Writing learners are provided with structured summaries, detailed notes, and text-based explanations generated using Gemini and BART models. Kinesthetic learners interact with virtual labs, simulations, drag-and-drop exercises, and embedded coding environments. This adaptive approach reduces cognitive overload while improving engagement and concept retention.
. 
\section{Retrieval-Augmented Generation Pipeline}
To maintain factual accuracy and syllabus alignment, the platform employs a Retrieval-Augmented Generation (RAG) pipeline. User queries are encoded using BERT embeddings and matched against a FAISS vector database containing curated educational content from engineering subjects. Relevant passages are retrieved based on cosine similarity and supplied to the Gemini language model for content generation. This process significantly reduces hallucination rates and ensures that all generated content remains curriculum aligned. The pipeline also supports modality-aware response generation, allowing the same concept to be delivered in multiple formats.

\section{Accessibility-First System Architecture}
Accessibility is integrated directly into the architecture of the platform instead of being treated as an optional feature. For visually impaired users, the system automatically prioritizes audio-based interaction and screen-reader compatibility. For hearing-impaired learners, captions, transcripts, and visual alternatives are provided for all audio content. Motor-impaired users are supported through keyboard navigation, gesture-based input systems, and simplified interaction pathways. The platform follows WCAG accessibility guidelines and allows users to manually override automatically inferred preferences whenever required.

\section{Gesture-Based Hardware Input Module}
The gesture input module is built using a fabric glove as the base, onto which flex sensors are affixed using adhesive along each finger. When the user forms a hand sign, the flex sensors bend and change their resistance proportionally. Each sensor is connected in a voltage divider configuration with a 10k$\Omega$ resistor, allowing the Arduino to read the analog voltage output corresponding to each finger’s position.

The Arduino reads these analog values in real time, processes them against a predefined set of gesture thresholds, and maps the combination to a corresponding character or word. The recognized output is simultaneously displayed on the LCD screen for immediate visual feedback to the user, confirming that the gesture was read correctly before it is passed to the platform.

The LCD display is connected to the Arduino via standard parallel or I2C interface and shows the detected gesture as readable text. This module enhances accessibility for users with speech or motor impairments and integrates directly into the adaptive learning pipeline. Once recognized, the gesture is converted into text and processed identically to typed or spoken input.

\section{Technical Stack and System Integration}
The frontend of the platform is developed using React.js, while the backend is implemented using Python Flask. The AI components include BERT for semantic embeddings, Gemini for text generation, BART for summarization, FAISS for vector retrieval, and Google Cloud Text-to-Speech for auditory content delivery. Data persistence is managed using PostgreSQL and MongoDB. The system follows a microservices architecture, ensuring scalability, modularity, and efficient deployment across institutional environments.